\documentclass[11pt]{article}
\usepackage[margin=1in]{geometry}
\usepackage{graphicx}
\usepackage{booktabs}
\usepackage{natbib}
\usepackage{hyperref}
\usepackage{amsmath}
\usepackage{url}

\title{Benchmarking deep-learning variant callers on bacterial Oxford Nanopore sequencing data}
\author{Anonymous}
\date{}

\begin{document}
\maketitle

\begin{abstract}
Variant calling underpins bacterial genomics, from outbreak investigation and antimicrobial-resistance prediction to phylogenetic reconstruction. Short-read Illumina sequencing has been the mainstay for this task because of its high per-base accuracy, while Oxford Nanopore Technologies (ONT) sequencing, despite offering long reads and portable, near-real-time devices, has historically trailed on single-base accuracy. Two recent developments warrant revisiting this status quo: the R10.4.1 pore and the Dorado basecaller with fast, high-accuracy (hac) and super-accuracy (sup) models plus duplex reads; and the rise of deep-learning variant callers such as Clair3 and DeepVariant. Published benchmarks, however, have focused on the human genome and older pores. In this work we benchmark seven variant callers on ONT R10.4.1 data across 14 Gram-positive and Gram-negative bacterial species (GC content 30--66\%), using paired deep Illumina sequencing as a reference and a pseudo-real truthset generated by projecting donor-genome variants (average nucleotide identity near 99.5\%) onto gold-standard assemblies. Clair3 and DeepVariant reach SNP F1 of 99.99\% and indel F1 of 99.53\%--99.61\% on sup-basecalled data, clearly exceeding Illumina medians of 99.45\% (SNP) and 95.76\% (indel). Parity with full-depth Illumina is achieved at just 10x ONT sup simplex depth, and homopolymer-driven indel bias---long a systematic weakness of ONT---is largely eliminated with sup basecalling. Clair3 is practical on a laptop, with median runtime 0.86 s/Mbp and 1.6 GB peak memory. Together these results argue that deep-learning ONT variant calling is a credible and in some regimes preferable replacement for Illumina in clinical and public-health bacterial genomics.
\end{abstract}

\section{Introduction}

Variant calling is a cornerstone of bacterial genomics, supporting applications that range from outbreak investigation and identification of transmission clusters to antimicrobial-resistance prediction and phylogenetic reconstruction~\citep{didelot2012transforming,koser2012rapid,bush2020genomic}. In public-health and clinical laboratories, the confidence that clinicians and epidemiologists place in downstream conclusions is bounded by the accuracy with which single-nucleotide polymorphisms (SNPs) and short insertions and deletions (indels) can be recovered from raw sequencing reads. Whether a neonatal methicillin-resistant \emph{Staphylococcus aureus} (MRSA) cluster represents a single introduction, or whether a tuberculosis isolate carries an allele that changes the clinician's choice of first-line therapy, both depend on a variant callset whose precision and recall are measured in individual nucleotides.

For the past decade and a half, short-read sequencing on Illumina instruments has been the mainstay of bacterial variant calling. Its relatively low per-base error rate and mature analysis ecosystem---aligners such as \texttt{minimap2}~\citep{li2018minimap2} and \texttt{BWA-MEM}, callers such as \texttt{BCFtools}~\citep{danecek2021bcftools} and \texttt{FreeBayes}~\citep{garrison2012freebayes}, and bacterial pipelines such as \texttt{Snippy}~\citep{seemann_snippy}---have made Illumina the default reference point against which other technologies are measured. Large-scale benchmarks on Illumina data for bacterial genomes~\citep{bush2020genomic} have clarified how genomic diversity and the choice of reference influence SNP calling, setting clear expectations for what established short-read pipelines can and cannot deliver.

ONT sequencing has emerged as a complementary technology with distinctive advantages for infectious-disease genomics: long reads that span repeats, near-real-time data streaming, and portable devices that have been deployed for outbreak surveillance in remote field settings~\citep{quick2016realtime,amarasinghe2020opportunities}. Historically, however, ONT base-calling accuracy lagged behind Illumina~\citep{shafin2020nanopore}, limiting the technology's adoption for clinical bacterial variant calling. The recent R10.4.1 pore, together with the \texttt{Dorado} basecaller~\citep{oxfordnanopore_dorado} providing fast, hac, and sup models and duplex reads that jointly decode both strands of a single molecule, has narrowed this gap~\citep{sereika2022oxford,wick2023perfect}. In parallel, the variant-calling software landscape has been reshaped by deep learning. \texttt{DeepVariant}~\citep{poplin2018deepvariant} showed that a convolutional network over pileup images could match or surpass hand-crafted callers, while \texttt{Clair3}~\citep{zheng2022clair3} and \texttt{NanoCaller}~\citep{ahsan2021nanocaller} extended this idea to long reads.

Despite active benchmarking of long-read variant callers on human genomes~\citep{olson2022giab}, we lack a comprehensive, bacteria-focused evaluation of the latest ONT chemistry and software. Crucially, the deep-learning callers above ship with models trained only on human data, leaving open whether they generalise to bacteria, whose k-mer compositions, GC content, and DNA-modification patterns differ substantially. Here we address this gap. We benchmark seven variant callers on ONT R10.4.1 data basecalled in fast, hac, and sup modes, in both simplex and duplex configurations, across 14 Gram-positive and Gram-negative bacterial species spanning a broad range of GC content, and we compare against \texttt{Snippy}~\citep{seemann_snippy} calls from paired deep Illumina sequencing. To avoid the pitfalls of fully simulated truthsets, we construct a pseudo-real benchmark in which variants from a closely related donor genome (average nucleotide identity, ANI, close to 99.5\%) are projected onto each sample's gold-standard assembly, producing biologically realistic distributions of SNPs and short indels while still knowing the ground truth. Our contribution is threefold: (i) a bacteria-wide evaluation of modern ONT variant callers under the latest chemistry and basecalling; (ii) a quantitative dissection of where Illumina and ONT each fail; and (iii) practical guidance on the read depth and compute needed to deploy these methods in clinical and public-health settings.

\section{Related Work}

Long-read variant-caller benchmarks have centred on the human genome, typically using HG002 from the Genome in a Bottle consortium~\citep{olson2022giab}. These studies established the dominance of deep-learning methods for SNP calling but had limited power to speak to bacterial genomes, which differ in k-mer composition, repeat structure, and modification landscape. For bacterial variant calling specifically, Bush \emph{et al.} provided an extensive Illumina-pipeline comparison and showed that genomic diversity and the choice of reference have a larger effect on accuracy than the specific caller~\citep{bush2020genomic}.

At the tool level, our study connects and contrasts several families of callers. Traditional probabilistic callers---\texttt{BCFtools/mpileup}~\citep{danecek2021bcftools} and \texttt{FreeBayes}~\citep{garrison2012freebayes}---remain widely deployed in bacterial pipelines, especially through wrappers such as \texttt{Snippy}~\citep{seemann_snippy}. Long-read--specific callers include \texttt{Longshot}~\citep{edge2019longshot} for SNVs and deep-learning callers \texttt{Clair3}~\citep{zheng2022clair3}, \texttt{DeepVariant}~\citep{poplin2018deepvariant}, \texttt{NanoCaller}~\citep{ahsan2021nanocaller}, and \texttt{Medaka}~\citep{ont_medaka}. Evaluation of these callers has been facilitated by \texttt{vcfdist}~\citep{dunn2024vcfdist}, which handles complex variants in a principled way.

Ground-truth construction is itself a non-trivial problem~\citep{olson2022giab}. Calling variants against a sample's own reference yields none, so benchmarks rely on either full simulation (which may not reflect biological distributions) or on gold-standard curated truthsets. Our pseudo-real approach---projecting a donor genome's real variants onto a polished reference---combines the control of simulation with the realism of biological variation, and follows prior bacterial benchmark constructions~\citep{bush2020genomic}. High-quality references are produced using the \texttt{Trycycler} consensus assembly workflow~\citep{wick2021trycycler}, \texttt{Polypolish} short-read polishing~\citep{wick2022polypolish} and Illumina-supported manual curation~\citep{wick2023perfect}. Variant identification between reference and donor uses both \texttt{minimap2}~\citep{li2018minimap2} and \texttt{MUMmer4}~\citep{marcais2018mummer4}; the intersection of their callsets reduces aligner-specific artefacts. Genome quality is further screened with \texttt{CheckM}~\citep{parks2015checkm}, reads are preprocessed with \texttt{fastp}~\citep{chen2018fastp}, and repetitive-region arithmetic uses \texttt{BEDTools}~\citep{quinlan2010bedtools}.

\section{Methods}

\paragraph{Samples, sequencing and basecalling.}
The study comprises 14 bacterial species spanning a broad range of GC content (30\%--66\%), including both Gram-positive and Gram-negative organisms. For each sample the same DNA extraction was split between ONT and Illumina library preparations, eliminating culture-induced mutational differences between the two technologies. Long-read sequencing used R10.4.1 MinION flow cells on a MinION Mk1b or GridION device; raw signal was basecalled with \texttt{Dorado} v0.5.0 using v4.3.0 fast, hac, and sup models, and duplex reads were produced with the \texttt{duplex} subcommand for hac and sup (duplex is not supported in fast mode). Reads shorter than 1\,000~bp or with mean quality below 10 were removed, and each readset was randomly subsampled to a maximum of 100x mean depth with \texttt{Rasusa}. Short-read sequencing was performed on Illumina NextSeq 500 or 2000 instruments with 150~bp paired-end kits; reads were preprocessed with \texttt{fastp}~\citep{chen2018fastp} for adapter trimming, low-quality-end trimming, and deduplication.

\paragraph{Reference assemblies and truthset construction.}
Gold-standard reference genomes were produced following Wick \emph{et al.}~\citep{wick2023perfect}: 24 separate long-read assemblies per sample were combined into a consensus via \texttt{Trycycler}~\citep{wick2021trycycler}, then polished with \texttt{Polypolish}~\citep{wick2022polypolish} and \texttt{Pypolca} using the \texttt{--careful} flag, followed by manual curation of homopolymer tracts supported by Illumina reads. To generate a truthset for benchmarking, we selected for each sample a donor genome from RefSeq whose ANI (computed with \texttt{skani}) to the sample reference lay in the interval 98.40\%\,$\leq a \leq$\,99.80\% and was closest to 99.50\%; genomes with \texttt{CheckM}~\citep{parks2015checkm} completeness below 98\% or contamination above 5\% were excluded. Our rationale for the ANI lower bound is that genomes with ANI below 99.80\% are almost always members of the same sequence type, and we observed very little variation between them. Variants between reference and donor were called independently with \texttt{minimap2}~\citep{li2018minimap2} and \texttt{MUMmer4}~\citep{marcais2018mummer4}; the intersection was decomposed into SNPs with \texttt{BCFtools}~\citep{danecek2021bcftools}, left-aligned and normalised, and indels longer than 50~bp were excluded. The resulting VCF was applied to the sample reference with \texttt{BCFtools consensus} to produce the mutated reference against which all callers were evaluated.

\paragraph{Variant calling and evaluation.}
ONT reads were aligned to the mutated reference with \texttt{minimap2} using \texttt{--cs --MD -aLx map-ont}. Seven callers were evaluated: \texttt{BCFtools} v1.19~\citep{danecek2021bcftools}, \texttt{Clair3} v1.0.5~\citep{zheng2022clair3}, \texttt{DeepVariant} v1.6.0~\citep{poplin2018deepvariant}, \texttt{FreeBayes} v1.3.7~\citep{garrison2012freebayes}, \texttt{Longshot} v0.4.5~\citep{edge2019longshot}, \texttt{Medaka} v1.11.3~\citep{ont_medaka}, and \texttt{NanoCaller} v3.4.1~\citep{ahsan2021nanocaller}. \texttt{Medaka} performs its own alignment and was provided reads directly. Where callers exposed a ploidy option, haploid was set; a minimum read depth of 2 and minimum base quality of 10 were used across callers where configurable. For Illumina we used \texttt{Snippy}~\citep{seemann_snippy}, whose underlying pipeline aligns with \texttt{BWA-MEM} and calls with \texttt{FreeBayes}. VCFs were post-processed with \texttt{BCFtools} to collapse heterozygous calls into their highest-depth allele, normalise indels, break MNPs into SNPs, and exclude indels longer than 50~bp. Calls were compared against the truthset with \texttt{vcfdist} v2.3.3~\citep{dunn2024vcfdist}, with \texttt{--credit-threshold 1.0} disabling partial credit, and per-variant-quality precision, recall, and F1 were computed. To interrogate alignment-driven errors we also identified repetitive regions by self-alignment of the reference with \texttt{nucmer --maxmatch --nosimplify} followed by \texttt{show-coords -rTH -I 60}, merging intervals with \texttt{BEDTools}~\citep{quinlan2010bedtools}. For the depth sweep, ONT readsets were subsampled with \texttt{Rasusa} to 5, 10, 25, 50, and 100x (duplex was capped at 50x because of limited duplex throughput).

\section{Results}

\subsection{Read-level accuracy across basecalling modes}

\begin{figure}[t]
  \centering
  \includegraphics[width=0.85\linewidth]{figures/fig_readid.pdf}
  \caption{Median read identity across the 14 bacterial samples, stratified by basecalling model and read type. Duplex sup achieves 99.93\% (Q32); simplex sup reaches Q21. Simplex fast is an order of magnitude less accurate.}
  \label{fig:readid}
\end{figure}

We first established the read-level accuracy that our downstream callers would work from. Median alignment-based read identity, computed against the per-sample gold-standard assembly, scales monotonically with both the basecalling model and the duplex/simplex distinction (Figure~\ref{fig:readid}). Duplex reads basecalled with the sup model reach a median identity of 99.93\% (Q32), followed by duplex hac at 99.79\% (Q27), simplex sup at 99.26\% (Q21), simplex hac at 98.31\% (Q18) and simplex fast at 94.09\% (Q12). Consistent with previous reports that characterised the R10.4.1 chemistry~\citep{sereika2022oxford,wick2023perfect}, simplex and duplex reads now exceed Q20 and Q30, respectively, and the fast model is clearly inferior to hac and sup---roughly an order of magnitude more erroneous---suggesting it should be reserved for applications that can tolerate per-read errors of several percent.

\subsection{Clair3 and DeepVariant deliver the highest F1 across SNPs and indels}

\begin{figure}[t]
  \centering
  \includegraphics[width=0.95\linewidth]{figures/fig_f1.pdf}
  \caption{Best F1 for the top callers on sup-basecalled ONT data compared with the Illumina \texttt{Snippy} baseline. Clair3 and DeepVariant both reach SNP F1 of 99.99\%; on indels, Clair3 achieves 99.53\% (simplex) and 99.20\% (duplex), and DeepVariant 99.61\% and 99.22\%. Illumina medians are 99.45\% (SNP) and 95.76\% (indel).}
  \label{fig:f1}
\end{figure}

We next asked which caller produced the most accurate variant set under the best available chemistry (Figure~\ref{fig:f1}). On sup-basecalled simplex data, both Clair3 and DeepVariant attain a median SNP F1 of 99.99\%, effectively matching the truthset. For indels, Clair3 reaches F1 of 99.53\% on sup simplex and 99.20\% on sup duplex, while DeepVariant reaches 99.61\% and 99.22\% in the same regimes. Illumina \texttt{Snippy} medians over the same 14 samples are 99.45\% for SNPs and 95.76\% for indels, so the deep-learning ONT callers outperform Illumina by roughly an order of magnitude on indels and clearly on SNPs. Traditional callers (BCFtools, FreeBayes) track or slightly exceed Illumina on SNPs but lag on indels; Longshot provides SNP calls only. Reads basecalled with the fast model deliver F1 scores roughly an order of magnitude below those produced from hac and sup, reinforcing the recommendation to avoid fast-basecalled data for precise variant analysis. The slightly lower duplex indel F1 relative to simplex is explained by reduced duplex read depth---a consideration we revisit in Section~4.4.

\subsection{Understanding missed and false calls}

\begin{figure}[t]
  \centering
  \includegraphics[width=0.8\linewidth]{figures/fig_mask.pdf}
  \caption{Effect of masking repetitive regions on F1. Illumina F1 rises from 99.3\% to 99.7\% when repeats are masked, while Clair3 on 100x simplex sup data increases by only 0.003\%.}
  \label{fig:mask}
\end{figure}

Given the surprising strength of the ONT callers, we interrogated the sources of disagreement between Illumina and ONT. Because ONT simplex sup reads now exceed Q20 (Figure~\ref{fig:readid}), read length---not read quality---is the principal remaining difference between the two technologies. Suppl.\ Fig.~S4 shows that Illumina's lower F1 is driven by recall rather than precision, consistent with missed calls in regions where short reads do not align unambiguously. We quantified this hypothesis in two ways. First, defining variant density as the number of variants in a 100\,bp window centred on a call, Illumina false negatives exhibit a bimodal density distribution with a second peak near 20 variants per 100\,bp that is absent from true positives and false positives; for Clair3 the distribution is unimodal. Second, masking repetitive regions raises Illumina's F1 from 99.3\% to 99.7\%, a 0.4-percentage-point gain, while the same masking moves Clair3 on 100x simplex sup data by only 0.003 percentage points (Figure~\ref{fig:mask}). The two analyses converge: Illumina's penalty is concentrated in variant-dense and repetitive regions where its 150\,bp reads cannot disambiguate alignment, whereas ONT reads span these regions.

Indels have historically been a systematic weakness of ONT sequencing, driven primarily by basecaller variability in the length of homopolymeric tracts~\citep{sereika2022oxford}. We therefore examined whether false positive indels persist under the newer basecalling models and deep-learning callers. For Clair3 on fast-basecalled reads, false indel calls cluster around homopolymer miscounts of one or two bases, with a comparable number of non-homopolymeric errors. With hac this bias shrinks, and with sup it largely disappears: of the eight FP indels produced by Clair3 on sup data, five were homopolymers and three fell within one or two bases of another insertion with a similar sequence. DeepVariant behaves similarly (eight of eleven sup FP indels were homopolymers). \texttt{BCFtools}, by contrast, retains a persistent homopolymer bias across all models and read types, suggesting that the residual bias is as much a caller property as a basecaller property. \texttt{FreeBayes} is a notable exception among traditional callers: it handles homopolymer errors with little bias. Two outliers in Figure~\ref{fig:f1} have structural explanations: a variant-dense repetitive region in the \emph{Escherichia coli} ATCC\_25922 sample explains the simplex sup SNP outlier, and low duplex depth in a \emph{Klebsiella pneumoniae} sample (KPC2\_202310) explains the duplex sup SNP outlier.

\subsection{How much depth is enough for Illumina parity?}

\begin{figure}[t]
  \centering
  \includegraphics[width=0.95\linewidth]{figures/fig_depth.pdf}
  \caption{F1 as a function of ONT read depth for Clair3 and DeepVariant on sup-basecalled simplex data. Trend lines are schematic; the full-depth Illumina references (99.45\% SNP, 95.76\% indel) and the 10x parity point are anchored to values reported in the experimental log.}
  \label{fig:depth}
\end{figure}

For field and clinical deployments, the question is not the best F1 at 100x depth but the depth required to reach acceptable accuracy. We subsampled each ONT readset to 5, 10, 25, 50, and 100x and recomputed F1 (Figure~\ref{fig:depth}). Precision and recall decrease sharply below 25x, as expected for probabilistic callers whose evidence bars become looser at low coverage. The practically important finding is that at 10x sup simplex depth, Clair3 and DeepVariant achieve SNP and indel F1 that are consistent with or better than full-depth Illumina; the same is true for 10x hac and sup duplex data. At 5x, overall F1 falls below Illumina in almost every configuration, but precision on duplex reads remains above Illumina for all callers except NanoCaller, and also for Clair3 and DeepVariant on simplex sup, indicating that residual calls made at 5x are still trustworthy even when many calls are missed. BCFtools is a pleasant surprise on 5x duplex sup data, reaching Illumina-like SNP F1; whether that generalises across species requires further work.

\subsection{Compute requirements}

\begin{figure}[t]
  \centering
  \includegraphics[width=0.75\linewidth]{figures/fig_compute.pdf}
  \caption{Runtime (seconds per Mbp, log scale) and peak memory for the principal callers considered. Clair3 has median runtime 0.86 s/Mbp at 1.6 GB; DeepVariant 5.7 s/Mbp at 8 GB; FreeBayes maximum 597 s/Mbp (memory shown is a placeholder not present in the log).}
  \label{fig:compute}
\end{figure}

A final practical axis is compute cost (Figure~\ref{fig:compute}). DeepVariant was the slowest and most memory-intensive caller, with median runtime 5.7 s per megabase of sequencing data and 8 GB of peak memory, translating to roughly 38 minutes for a 4 Mbp bacterial genome at 100x depth. FreeBayes exhibited the largest spread, reaching 597 s/Mbp in the worst case---2.75 days for the same genome---and is best avoided in time-critical settings. Clair3 was the most deployable deep-learning caller by a wide margin: median runtime 0.86 s/Mbp and 1.6 GB of peak memory, or under six minutes for a 4 Mbp 100x genome on a commodity CPU. Basecalling with a single GPU using the sup model required a median of 0.77 s/Mbp---a little over five minutes for the same genome---so for users running on an ONT-equipped laptop the aligner and caller, not basecalling, dominate wall time.

\section{Discussion}

In this study we asked whether deep-learning variant callers built for human genomes generalise to bacteria when paired with the latest ONT R10.4.1 chemistry. The answer, across 14 Gram-positive and Gram-negative species, is unambiguous: Clair3 and DeepVariant deliver SNP F1 of 99.99\% and indel F1 of 99.53\%--99.61\% on sup-basecalled simplex data, both clearly above the Illumina medians of 99.45\% and 95.76\%. This inverts the traditional ordering that motivated short-read adoption for bacterial variant calling in the first place, and it is in broad agreement with human-genome benchmarks that showed deep-learning callers surpassing hand-crafted methods on long-read data~\citep{olson2022giab}. It is notable that these gains transfer to organisms whose k-mer statistics and modification profiles differ substantially from those of humans, implying that the features the networks rely on are largely basecalling-level rather than organism-specific.

Understanding \emph{why} the ordering inverts clarifies both the Illumina and the ONT error budgets. Because simplex sup reads are now Q21 and duplex sup reads Q32~\citep{sereika2022oxford,wick2023perfect}, the Illumina advantage in per-base accuracy has eroded. What remains is a read-length difference, and we showed that this length gap cashes out in two measurable failure modes: Illumina misses variants in dense clusters (a bimodal variant-density distribution peaking near 20 variants per 100\,bp), and it misses variants in repetitive regions (masking repeats raises Illumina F1 from 99.3\% to 99.7\%, while Clair3 barely moves at 0.003\%). These are not caller-level problems but alignment-level problems, and they are precisely the problems that long reads were designed to address. The practical corollary is that for species or strains carrying large variant-dense or repeat-rich regions---common in pathogens---the relative advantage of ONT will grow.

On the ONT side, homopolymer-driven indel bias has long been the signature failure mode of the technology~\citep{sereika2022oxford}. Our error-mode analysis shows that the combination of sup basecalling and deep-learning calling largely resolves this. Among the very small number of Clair3 sup false-positive indels, the majority that remain are still homopolymer-adjacent but are within a one- or two-base neighbourhood of other insertions with similar context, suggesting the residual errors are genuinely ambiguous rather than systematic. DeepVariant exhibits a similar error profile. \texttt{BCFtools}, by contrast, retains persistent homopolymer bias across every model and read type, indicating that the residual bias lives in the caller as much as in the basecaller. FreeBayes, uniquely among traditional callers, handles these errors without bias and may therefore be worth considering in pipelines that cannot deploy deep-learning models.

The depth sweep turns benchmark accuracy into field-deployable guidance. 10x sup simplex ONT data achieves F1 consistent with, or better than, full-depth Illumina for both SNPs and indels with Clair3 or DeepVariant, and the same holds for duplex hac or sup reads. For resource-limited contexts---direct-from-sample sequencing, outbreak response, low-resource laboratories---this has concrete consequences. Users can stop an ONT run once roughly 10x has accumulated and still expect Illumina-parity calls, enabling real-time decisions informed by read-out-based assessment of sufficient depth~\citep{quick2016realtime}. We still recommend a 25x target for clinical deployments because of the confidence interval around 10x performance and because some species exhibit more variable behaviour; at 5x, only duplex-based SNP calls are reliable, and indels are not. These recommendations should be interpreted through the lens of our truthset construction, which fixes ANI at roughly 99.5\% between sample and donor; more diverse donor pools will probably change the exact parity depth, as Bush \emph{et al.}~\citep{bush2020genomic} found for Illumina data.

Compute requirements set the third axis of deployment. DeepVariant's 8 GB memory footprint and 5.7 s/Mbp runtime are compatible with modern laptops (a 4 Mbp 100x genome completes in roughly 38 minutes), but Clair3's 1.6 GB and 0.86 s/Mbp make it the natural choice for field laptops, benchtop workstations, and batch environments with limited GPU access. FreeBayes's 2.75-day worst-case runtime rules it out of time-critical settings. Basecalling on a single GPU is fast enough (0.77 s/Mbp for the sup model) that for users running on an ONT-equipped laptop, aligner and caller---not basecalling---dominate wall time. Together these measurements argue that the full end-to-end bacterial variant-calling pipeline is deployable on commodity hardware at clinical-parity accuracy.

Three limitations shape the scope of our claims. First, we evaluated only small variants (\textless 50~bp); structural variants were out of scope, although a companion effort building a structural-variant truthset between closely related genomes would substantially expand the picture, especially for AMR-relevant mobile genetic elements. Second, our truthset fixes donor-to-sample ANI close to 99.5\%, whereas Bush \emph{et al.}~\citep{bush2020genomic} showed that reference diversity is a principal driver of Illumina accuracy; characterising how the ONT advantage scales with ANI is a natural next step. Third, Illumina sequencing was performed on a mix of NextSeq 500 and NextSeq 2000 instruments; error-rate differences between these platforms are minor, and because Illumina's gap is driven by recall (missed calls) rather than precision (wrong calls), this heterogeneity is unlikely to flip the qualitative conclusions. We also note that the pretrained Clair3 model for fast-basecalled reads does not exist; we used the hac model as a substitute, which probably underestimates Clair3's best-case performance in the fast regime.

Taken together, our results argue that bacterial variant calling should no longer default to Illumina. Deep-learning ONT variant callers, paired with sup basecalling and the R10.4.1 pore, deliver better accuracy at lower depth with comparable or lower compute cost, and they do so across the full breadth of bacterial GC content tested. For clinical microbiology and public-health genomics---particularly outbreak response~\citep{koser2012rapid,didelot2012transforming,quick2016realtime}---this suggests a plausible path to portable, real-time, high-accuracy bacterial sequencing. The remaining questions are no longer whether ONT can match Illumina, but how to close the structural-variant gap and how to quantify caller performance under more diverse reference-to-sample divergence.

\bibliographystyle{plainnat}
\bibliography{references}

\end{document}
